% HD-G33. GPT-6 Astra Pro /root/g34; inherited effort not exposed.
% 21 November 2026. Actual finite-group and representation closure.
\section{HD-G33: closure of the primitive rank-three branch}
\label{sec:primitive-rank-three-finite-closure}

Retain the actual natural sixfold source of HD-G31--G32, with
\(K=\mathbb Q(i)\) or
\(K=\mathbb Q(\sqrt{-d})\), \(d=2,7,11,19,43,67,163\), exact
coefficient character field \(K\), period rank nine, and no
nonidentity coefficient-group element of moving rank one or two.
The normal subgroup \(N_3\) generated by its actual moving-rank-three
elements cannot be absolutely irreducible and complex primitive.
This closes the finite range left by HD-G31--G32. The proof keeps
the original finite group, its actual normalizing matrices and its
primitive Jacobian factor throughout. In particular it does not
replace an actual outer action by a putative extension character.

\paragraph{1. Conjugate generation inside the original matrix group.}
HD-G31 and HD-R14 give \(15\le m\le35\) in this branch. HD-G32
gives a characteristic absolutely irreducible quasi-simple subgroup
\(S\triangleleft N_3\), normalized by the original \(G\), and
the actual injection
\[
 G/Q\hookrightarrow\operatorname{Aut}(L),\qquad
 L=S/Z(S),\quad Q=G\cap K^*I,
 \qquad |Z(S)|\mid\gcd(m,|\mu(K)|).                 \tag{G33.1}
\]
There is an actual involution \(y\in G\) of moving rank three.
Indeed HD-G31 gives either such a generator already, or the mixed
Gaussian spectrum \((i,i,-i)\) or its conjugate; its actual square
has spectrum \((-1,-1,-1)\) on its moving space. The image
\(\bar y\) in (G33.1) is a nonidentity involution, since a scalar
in degree at least15 cannot have moving rank three.

For a nonidentity automorphism \(x\) of a simple group \(L\),
write \(\alpha(x)\) for the least number of \(L\)-conjugates of
\(x\) generating \(\langle L,x\rangle\). This is the exact
definition in Guralnick--Saxl, \emph{Generation of finite almost
simple groups by conjugates}, J. Algebra268 (2003),519--571,
p.519. Lift the conjugating elements to \(S\), and let \(J\)
be generated by the resulting actual conjugates \(y_i\) of \(y\).
For every \(s\in S\) there is \(j\in J\) with the same image in
\(\operatorname{Aut}(L)\), so \(S\subset QJ\). Perfection gives
\[
 S=[S,S]\subset[QJ,QJ]=J'.
\]
The common fixed space of all \(y_i\) is zero, since \(S\) is
nontrivial and irreducible. Intersection of their fixed spaces
therefore gives the useful bound in the original representation:
\[
                        m\le3\alpha(\bar y).                \tag{G33.2}
\]
Every lift here belongs to the original \(G\); no scalar has been
adjoined to it.

\paragraph{2. The complete corrected finite filter.}
We use Hiss--Malle, LMS J. Comput. Math.4 (2001),22--63,
Theorem1 p.22 and Section6 p.30, together with the replacement
table in their \emph{Corrigenda}, volume5 (2002),95--126.
The correction instructions are on pp.95--96; all entries of
degrees15--35 are on pp.99--104. The generic alternating and
\(L_2(q)\) families in the original Table2 p.31 remain in force.
Characteristic zero is retained also in a table entry denoted
\(\ell\ne p_1,p_2,\ldots\). Covers whose centres have order
divisible by three are excluded by (G33.1). Other field and centre
restrictions may be ignored in the following larger list, so that
the filter cannot discard a representation through an unproved
field identification.

Here is the resulting list of simple quotients, possible degrees,
and a sufficient uniform bound \(a\ge\alpha(\bar y)\).
Degrees coming from different permissible covers are combined.
The symbol \(S_{2r}(q)\) in Hiss--Malle means the simple
symplectic group, and is written \(\operatorname{Sp}_{2r}(q)\)
below when \(q\) is even.
\[
\begin{array}{c|l|c}
 L&\hbox{degrees in the larger list}&a\\\hline
 A_7&15,20,21,35&6\\
 A_8&20,21,24,28,35&7\\
 A_9&21,27,28,35&8\\
 A_{10},A_{11}&16,35\ ;\ 16&9\ ;\ 10\\
 A_{12},A_{13}&32\ ;\ 32&11\ ;\ 12\\
 L_3(3),L_3(4),L_3(5)&16,26,27\ ;\ 20,28,35\ ;\ 30,31&4\\
 L_4(3),L_5(2)&26\ ;\ 30&6\ ;\ 5\\
 U_3(3),U_3(5)&21,27,28,32\ ;\ 20,21,28&4\\
 U_4(2),U_4(3)&15,20,24,30\ ;\ 20,21,35&6\\
 U_6(2)&22&6\\
 \operatorname{Sp}_4(4),\operatorname{PSp}_4(7)&18,34\ ;\ 24,25&5\\
 \operatorname{Sp}_6(2),\operatorname{Sp}_8(2)&15,21,27,35\ ;\ 35&7\ ;\ 9\\
 O_8^+(2),O_8^-(2)&28,35\ ;\ 34&8\\
 {}^3D_4(2),{}^2F_4(2)',\operatorname{Sz}(8)&26\ ;\ 26,27\ ;\ 35&7\ ;\ 8\ ;\ 5
\end{array}                                                    \tag{G33.3}
\]
The remaining sporadic quotients are
\[
\begin{array}{c|l}
L&\hbox{degrees}\\\hline
M_{11},M_{12},M_{22},M_{23},M_{24}&16\ ;\ 16,32\ ;\ 21\ ;\ 22\ ;\ 23\\
J_2,HS,McL,Co_3,Co_2,Co_1,Ru&21\ ;\ 22\ ;\ 22\ ;\ 23\ ;\ 23\ ;\ 24\ ;\ 28.
\end{array}                                                    \tag{G33.4}
\]
For each quotient in (G33.4), \(\alpha(\bar y)\le5\).

The bounds just used are the complete applicable conclusions of
Guralnick--Saxl Lemma3.1 pp.527--529, Theorems4.1--4.4
pp.533--534, Theorem5.1 p.541 (with Propositions5.5,5.7,5.8
pp.545--547), Lemma6.1 pp.547--548, and Table1 p.554 with
the calculation in the proof of Lemma7.6 pp.553--554.
In particular \(L_2(q)\), \(q\ne9\), has \(\alpha\le4\) for
an involution in its full automorphism group. The exceptional
\(L_2(9)=A_6\) has no retained degree15--35 row: its degree15
row has centre divisible by three. Generic \(A_n\), \(n\ge14\),
has only the ordinary augmentation degree \(n-1\) in this range.
An actual element of \(S\) mapping to a three-cycle has moving
rank two in that literal augmentation model; the finite
\(K\)-intertwiner of HD-G31 transports this element back to the
original group. Thus the generic alternating family is excluded.

Combining the remaining entries with (G33.2) leaves exactly
\[
\begin{split}
 &(A_7,15),\ (A_8,20),\ (A_8,21),\ (A_9,21),\\
 &(2.A_{10},16),\ (2.A_{11},16),\ (2.A_{12},32),\ (2.A_{13},32),\\
 &(U_4(2),15),\ (\operatorname{Sp}_6(2),15),\
                   (\operatorname{Sp}_6(2),21).
\end{split}                                                    \tag{G33.5}
\]
For the last three rows the centre is trivial by (G33.1), since
the degree is odd. The corrected tables are used only after the
actual quasi-simple reduction, and every remaining normalizer
case is treated below.

\paragraph{3. Literal Coxeter lifts for all alternating rows.}
For \(A_n\), \(n\ge7\), the cited Lemma6.1 sharpens the bound
to \(\alpha(x)\le\lfloor n/2\rfloor\) unless \(x\) is a
transposition. In each alternating case in (G33.5), (G33.2)
therefore forces \(\bar y\) to be a transposition in
\(\operatorname{Aut}(A_n)=S_n\).
Choose actual \(S\)-conjugates \(T_1,\ldots,T_{n-1}\) of \(y\)
whose images are the adjacent transpositions. Any two of their
fixed spaces intersect in dimension at least \(m-6>0\).
For \(|i-j|>1\), \((T_iT_j)^2\in Q\); for \(j=i+1\),
\((T_iT_j)^3\in Q\). Each scalar fixes this nonzero intersection
pointwise, and hence equals one. Along with \(T_i^2=1\), these
are the exact Coxeter relations. They define an actual map
\[
 \Phi:S_n\longrightarrow G,\qquad (i,i+1)\longmapsto T_i,
 \qquad \overline\Phi=\operatorname{id}_{S_n}.              \tag{G33.6}
\]
Thus \(\Phi\) is injective. Moreover \(S\subset Q\Phi(S_n)\),
so perfection gives \(S\subset\Phi(A_n)\). Their images in
\(A_n\) coincide, and \(\Phi(A_n)\) has trivial scalar kernel;
consequently \(S=\Phi(A_n)\). This excludes every double-cover
row in (G33.5) without assuming that its projective action lifts
to the usual symmetric-group representation.

The actual remaining \(S_n\)-module is irreducible, because its
restriction to \(S\) is. The ordinary Young classification and
hook formula give the following complete small-degree lists:
\[
\begin{array}{c|c|l}
n&m&\hbox{partitions}\\\hline
7&15&(5,1,1),\ (3,1,1,1,1)\\
8&20&(6,2),\ (2,2,1,1,1,1)\\
8&21&(6,1,1),\ (3,1,1,1,1,1)\\
9&21&\varnothing.
\end{array}                                                    \tag{G33.7}
\]
One primary account is Etingof et al., \emph{Introduction to
representation theory}, Student Mathematical Library59 (2011),
the author version dated10 January2011, Theorem4.36 and proof
pp.58--60, and Theorem4.53 and proof pp.63--64. In using its
Young construction we use unnormalized row and column sums;
their product \(c\) satisfies
\(c^2=(n!/\dim V_\lambda)c\). The averaged row--column product
is exactly \(c/(|R_\lambda||C_\lambda|)\), so its square is
\(n!/(\dim V_\lambda|R_\lambda||C_\lambda|)\) times itself.
These literal scalar maps fix the convention independently of
the averaged notation in that account.

The models and their relevant maps can also be given directly.
Let \(W=\ker(\mathbb Q^n\xrightarrow{\sum}\mathbb Q)\).
For \(n=7,8\), \(\Lambda^2W\) has dimensions15,21 and
transposition traces5,9, respectively; its sign twists have the
opposite traces. For the degree20 case let \(M\) have basis the
two-subsets of eight labels and set
\[
 J:\mathbb Q^8\longrightarrow M,\qquad
 J(e_i)=\sum_{j\ne i}e_{\{i,j\}},\qquad
 J^tJ=6I+\boldsymbol1\boldsymbol1^t.
                                                               \tag{G33.8}
\]
The rational invariant complement \(\ker J^t\) has dimension20
and transposition trace10, and its sign twist has trace\(-10\).
It is irreducible: the action on pairs has three orbitals, while
\(J(\mathbb Q^8)=1\oplus W\), leaving character norm one.
For \(\Lambda^2W\), at a permutation with \(f\) fixed labels
and \(t\) two-cycles, the character is
\(((f-1)^2-(f+2t-1))/2\); summing its square over the cycle
partitions with weights \(n!/\prod_j j^{a_j}a_j!\) gives norm
one for both displayed \(n\). The exact finite hook enumeration
and these norm sums are included in the certificate. Thus the
distinct model/sign-twist pairs exhaust (G33.7).

An involution of moving rank three in dimension \(m\) has trace
\(m-6\). None of the displayed traces is \(9,14,15\),
respectively. The actual \(K\)-linear identification is obtained
by solving
\(C\rho(T_i)=R((i,i+1))C\) for these rational models. The
complex isomorphism gives a one-dimensional solution space;
the coefficient equations have the same rank over \(K\), so
a nonzero \(K\)-solution exists and is invertible. This closes
the alternating cases inside the original representation.

\paragraph{4. A constructive exact character calculation.}
For the two remaining simple groups we use actual rational root
and coset matrices, rather than a stored irreducible-character
table. The common calculation is as follows. Enumerate the
conjugacy classes of the actual permutation group, their sizes,
representatives and power maps for powers one through five.
Start with the rational representations described below.
Form exterior powers of degrees two through five, symmetric
squares and cubes, and tensor products. Newton's identities
compute their character values exactly. For each known rational
irreducible \(V_j\), its full isotypic part in an actual module
\(M\) is the image of the literal rational projector
\[
 e_j(M)=\frac{\dim V_j}{|L_0|}
           \sum_{g\in L_0}\chi_j(g^{-1})\,M(g).
                                                               \tag{G33.9}
\]
All characters constructed here are integer valued. Subtracting
all these mutually orthogonal projectors therefore leaves an
actual rational summand. Its character has norm one exactly when
that nonzero summand is absolutely irreducible. This defines
every recorded complement by an actual map, including its
inclusion and projection. The algorithm retains only such
norm-one residuals and records each operation and each removed
nonnegative multiplicity.

For the first root group this procedure constructs25 pairwise
orthogonal rows, whose squared degrees sum to51840, the group
order, and there are25 conjugacy classes. For the second it
constructs30 such rows, with squared-degree sum1451520 and30
classes. These equalities prove completeness of the actual
irreducibles constructed. No unconstructed character is used to
assert a representation exists. The exact code uses permutation
group order, class enumeration and coset enumeration algorithms;
this is a reproducible finite arithmetic computation, not a claim
of verification in a formal proof kernel.

\paragraph{5. The actual \(U_4(2)\) group and its whole normalizer.}
Let \(C\) be the integral Cartan matrix with diagonal2 and
entries\(-1\) at the edges
\((1,3),(3,4),(4,5),(5,6),(2,4)\).
On the72-vector orbit of the coordinate roots use the column
reflections \(s_i=I-e_iC_{i,*}\). The resulting permutation
group \(W\) has order51840 and centre one; its derived subgroup
\(H\) has order25920 and is simple. These statements are checked
from these permutations. To avoid an action-convention shortcut,
permutation products act on the right, and their rational matrix
is the transpose of the corresponding column action. In the
basis of the coordinate roots it is recovered by taking as rows
the images of these roots. This is a genuine homomorphism,
verified on all generator pairs.

Here is an explicit group identification. The recorded matrices
\(a,b\in H\) generate \(H\) and satisfy
\[
 a^2=b^5=(ab)^9=[a,b]^3=[a,bab]^2=1.                         \tag{G33.10}
\]
Todd--Coxeter enumeration for this exact presentation gives
index12960 for the subgroup generated by \(a\). Thus its
presented group has order at most25920; the displayed image has
order25920, so the presentation is complete. In the standard
four-dimensional unitary group over
\(\mathbb F_4=\mathbb F_2[z]/(z^2+z+1)\), the actual pair
\[
 a_u=\begin{pmatrix}1&0&0&1\\0&1&0&0\\0&0&1&0\\0&0&0&1\end{pmatrix},
 \qquad
 b_u=\begin{pmatrix}
 0&0&1&1\\0&0&z&z^2\\1&1&z^2&z\\z^2&z&z^2&z
 \end{pmatrix}                                             \tag{G33.11}
\]
preserves the anti-diagonal unitary form, has determinant one,
satisfies (G33.10), and generates25920 matrices. The unitary
order formula is
\(|SU_4(2)|=2^6(2^2-1)(2^3+1)(2^4-1)=25920\), with
trivial centre. Hence the generator map
\(a_u\mapsto a,b_u\mapsto b\) is an actual isomorphism
\(SU_4(2)\xrightarrow{\sim}H\). This is the group denoted
\(U_4(2)\) in (G33.5).

We also compute its full automorphism group, so that outer
normalizers are included. The two involution classes of \(H\)
have different sizes45 and270, and the chosen \(a\) lies in
the class of size45. There are exactly1152 elements \(b'\)
of order five satisfying (G33.10) with \(a\) and with
\(|ab'|=9\). Each pair defines an automorphism: it defines a
nontrivial homomorphism from the complete presentation of the
simple \(H\), and therefore an isomorphism onto \(H\).
Conversely every automorphism fixing \(a\) arises this way.
Every automorphism preserves its uniquely sized involution
class, and inner automorphisms are transitive on it. Consequently
\[
                |\operatorname{Aut}(H)|=45\cdot1152=51840.
                                                               \tag{G33.12}
\]
The six-dimensional root representation restricts irreducibly
to \(H\), as its exact character norm shows. The centralizer
of \(H\) in \(W\) is therefore scalar, hence central in \(W\)
and trivial. Conjugation gives an injection
\(W\to\operatorname{Aut}(H)\); (G33.12) makes it an isomorphism.
This proves the full actual automorphism identification needed
for \(y\), without assuming that its action is inner.

Use the four seeds \(1,\det,V_6,\det\otimes V_6\) in
(G33.9). The complete construction gives four degree15 rows.
In the order of the involution class sizes
\(540,36,45,270\), their traces are
\[
 \begin{pmatrix}
 -3&5&-1&-1\\3&-5&-1&-1\\1&5&7&3\\-1&-5&7&3
 \end{pmatrix}.                                             \tag{G33.13}
\]
None of these entries is \(9\) or \(-9\). In addition, every
degree15 or degree30 irreducible of \(W\) constructed above
has restriction norm one on the actual index-two subgroup
\(H=\ker\det\).

This last check treats the inner case without assuming an
extension. Every irreducible \(H\)-module of dimension15 occurs
in the restriction of an irreducible summand of its actual
induced \(W\)-module. Index-two restriction has either one
irreducible constituent or two conjugate constituents, each
with multiplicity one: the isotypic decomposition and an outer
coset representative give the two alternatives directly; in
the invariant case choose its intertwiner and rescale its square
over \(\mathbb C\). The possible \(W\)-degrees are15 and30.
All degree30 restrictions are irreducible of dimension30, by
the recorded norm, and cannot contain that15-dimensional
module. Thus every relevant \(H\)-module is the restriction of
one of the four actual rational degree15 models in (G33.13).

Identify the original centre-free \(S\) with this \(H\).
When \(y\) induces an inner automorphism, \(y=\lambda s\)
for \(s\in S\). Its involutory image gives \(s^2=1\), and
\(y^2=1\) gives \(\lambda=\pm1\). Formula (G33.13) excludes
trace9. When its image is outer, the actual group
\(\langle S,y\rangle=S\rtimes\langle y\rangle\) has order51840
and maps faithfully by conjugation onto \(\operatorname{Aut}(S)\).
It is therefore the actual \(W\) just identified, and its
irreducible degree15 representation is one of (G33.13).
Again its involution cannot have trace9. The rational
projector models and the original matrices are related by the
finite \(K\)-intertwiner equations, whose nonzero solutions are
invertible by irreducibility, exactly as in paragraph3.

\paragraph{6. The actual \(\operatorname{Sp}_6(2)\) group.}
In the sum-zero hyperplane of \(\mathbb Q^8\), let the126
scaled roots be
\[
 2(e_i-e_j)\quad(i\ne j),\qquad
 (\epsilon_1,\ldots,\epsilon_8),\quad
       \epsilon_i\in\{1,-1\},\quad\sum_i\epsilon_i=0.
                                                               \tag{G33.14}
\]
Use the seven negative adjacent coordinate swaps and the map
\[
             v\longmapsto-v+\frac{v\cdot h}{4}h,
       \qquad h=(1,1,1,1,-1,-1,-1,-1).                      \tag{G33.15}
\]
They are determinant-one permutations of this root set. Their
group \(L_0\) has order1451520; the subgroup \(D\) generated
by the first seven is the literal signed permutation group
\(\{\operatorname{sgn}(\pi)P_\pi:\pi\in S_8\}\), of order40320.
Thus \([L_0:D]=36\). The root matrices are recovered from the
images of \(2(e_i-e_8)\), with the same right-action convention
as above; all generator products agree.

Here is the exact finite-field bridge. Choose the seven lattice
vectors
\[
 2(e_i-e_{i+1})\ (1\le i\le6),\qquad
             (-1,-1,-1,1,1,1,1,-1).                        \tag{G33.16}
\]
Every root in (G33.14) has integral coordinates in this basis.
The Gram matrix, with dot products divided by four, has
determinant two. Modulo two its radical is one-dimensional;
the quotient \(E\) is a six-dimensional symplectic space.
The63 pairs \(\{r,-r\}\) map bijectively to \(E\setminus\{0\}\).
The eight generators induce literal transvections
\(T_v(x)=x+B(x,v)v\). Their permutation group on these63
vectors has order
\[
         2^9(2^2-1)(2^4-1)(2^6-1)=1451520
                    =|\operatorname{Sp}_6(2)|.             \tag{G33.17}
\]
Hence the actual homomorphism \(L_0\to\operatorname{Sp}(E)\)
is an isomorphism. All integral coordinates, the radical,
the generator transvection identities and both permutation
orders are reconstructed independently in the supplementary
Python implementation.

The full automorphism group is inner, by the following argument
on this actual model. Its involution classes have sizes
\(63,315,945,3780\). Thus the transvection class is preserved
by every automorphism. Such an automorphism gives a bijection
\(f:E\setminus\{0\}\to E\setminus\{0\}\) preserving commuting
transvections, equivalently preserving the value of \(B\).
Put \(f(0)=0\). For every \(u,v,w\), preservation of \(B\)
and its bilinearity give
\[
 B(f(u+v),f(w))=B(u+v,w)
      =B(f(u)+f(v),f(w)).
\]
Surjectivity and nondegeneracy give additivity, so
\(f\in\operatorname{Sp}(E)\). Conjugation by \(f\) has the
same action on every transvection. They generate the group,
as is already witnessed by (G33.15)--(G33.17), so the original
automorphism is this inner automorphism. This proves the full
normalizer assertion used below.

\paragraph{7. Actual signed cosets and all degree15 and21 modules.}
Let \(\epsilon:D\to\{\pm1\}\) be the sign of its underlying
permutation, and use the actual induced module
\(M=\operatorname{Ind}_D^{L_0}\epsilon\). For recorded right
coset representatives \(t_i\), its row matrix is explicitly
\[
 M(g)_{ij}=\epsilon(t_i g t_j^{-1})
       \quad\hbox{when }Dt_i g=Dt_j,
       \qquad M(g)_{ij}=0\hbox{ otherwise}.                \tag{G33.18}
\]
The formula verifies \(M(gh)=M(g)M(h)\) directly by multiplying
the two subgroup elements. In the actual signed36 module,
let \(A\) be the sum of the63 transvection matrices. It is a
central self-adjoint operator and satisfies
\[
 (A+21I)(A+33I)=0,\qquad
 P_{15}=\frac{A+33I}{12},\quad
 P_{21}=-\frac{A+21I}{12}.                                  \tag{G33.19}
\]
These are complementary rational orthogonal projectors of
ranks15 and21. Their characters each have norm one. Along
with \(1,V_7,\Lambda^2V_7\) they are the five seeds for
(G33.9). The complete30-row construction shows that the only
degree15 module is \(\operatorname{im}P_{15}\), and the only
degree21 modules are \(\Lambda^2V_7\) and
\(\operatorname{im}P_{21}\). In the order of involution class
sizes \(3780,63,945,315\), the three character rows are
\[
\begin{array}{c|rrrr}
15&-1&-5&3&7\\
21&-3&9&1&-3\\
21&-3&-11&5&5.
\end{array}                                                   \tag{G33.20}
\]
Neither sign of an entry equals \(m-6\) in the respective row.

The original centre-free \(S\) is now identified with this
actual \(L_0\). By the full automorphism argument, every
original normalizer has the form \(\lambda s\) with \(s\in S\).
For our \(y\), \(s^2=1\) and \(\lambda^2=1\). Each actual
irreducible \(K\)-model of degree15 or21 is carried to one of
the rational modules above by the finite \(K\)-intertwiner
equations. Thus (G33.20) contradicts \(\operatorname{tr}y=m-6\)
in all remaining cases.

\paragraph{8. Original geometric factor and verification scope.}
The contradiction has excluded the complete primitive branch,
including every original scalar correlation and outer action.
All elements \(y,y_i,T_i\), the subgroup \(S\) and its scalar
kernel are actual elements and subgroups of the original
coefficient group. The finite model maps in paragraphs3,5--7
are group homomorphisms and invertible \(K\)-intertwiners;
their representations are actual tensor, exterior, coset and
projector summands. For any such coordinate map \(C\), retain
the original primitive line as \(KCw\) and the transported
positive Hermitian form. A prescribed invariant model form
instead satisfies \(h_0(Cx,Cy)=a h(x,y)\), \(a\in\mathbb Q_{>0}\),
and the normalized coefficient
\[
 \frac{h_0(Cw,C\rho(g)w)}{h_0(Cw,Cw)}
                    =\frac{h(w,\rho(g)w)}{h(w,w)}           \tag{G33.21}
\]
is identical. Therefore the rational group-algebra idempotent,
its integer-cleared image in the same covering Jacobian, its
restricted theta class and its corner \(K\)-action are unchanged.
The two Schur factors are \(c_0=c/a\), so \(c_0a=c\).
There is no unproved geometric replacement in the finite filter.

The exact GAP program in \path{rank_three_review/},
\path{construct_small_core_characters.g},
constructs the actual groups, presentation bridge, automorphism
count, signed-coset projectors and every irreducible recipe.
It emits both complete class data sets and all recipe operands.
Its independent rational-model reconstruction is
\path{rank_three_review/independent_small_core_models.py};
the independent field, Young-model and finite-filter checks are
\path{rank_three_review/primitive_core_filter.py}.
A separate reviewer implementation replays all Newton operations,
all removed multiplicities, orthogonality and completeness from
the emitted class enumeration, power maps and seed characters.
That replay does not independently enumerate classes or seeds.
The second model implementation reconstructs the seeds, root
groups, mod-two bridge and signed36 projectors; it accepts the
recorded class representatives and the separately checked
Todd--Coxeter presentation bound. These scopes are recorded
separately in the check ledger.
